% Transcription — fonds Alexandre Grothendieck, shelfmark 76, batch 1 (pages 1-20).
% Established from the facsimile archives/batches/76/batch-01.pdf, cut from a
% copy of 76.pdf fetched through the project's relay
% (https://grothendieck-relay.onrender.com/source/76.pdf); PDF page k =
% archivists' page k, no cover sheet. Per the skill transcribe-grothendieck.
% The folder is sixty-two pages in four batches (1-20, 21-40, 41-60, 61-62),
% transcribed in parallel by four passes; this is the first, and the others'
% files could not be relied on.
% Pass: Opus 5.5 (claude-opus-5-5), 2026-09-26 — first pass, unchecked against the pages by a human.
% Caveat: pages 15-20 were checked against the facsimile only at whole-sheet
% resolution, not tile by tile; they are the first to verify.
%
% Dating and title. The dating is the folder's own, copied verbatim from the
% catalogue; the group « Géométrie et topologie combinatoire » (69-89) holds
% it, and since the folder has a date of its own the group's range is not
% added. The title is the catalogue's, without its final full stop.
%
% Pages skipped: 1, 10 and 14, blank sheets (1 and 14 cream, with only
% show-through; 10 blank). Pages summarised: none; nothing is typed and
% nothing is by another hand.
%
% Pages transcribed from his hand: 2-9, 11-13, 15-20. Pp. 2-9 are in
% blue-black ink on perforated listing paper, some of it ruled like music
% paper (pp. 3, 6, 8 are sideways sheets); a few words on pp. 2-3 are added
% in black ink. Pp. 11-13 and 15-19 are bright blue felt-tip, p. 20 black
% felt-tip. The hand is fast drafting: formulas carry, much connective
% prose does not and is left \ill{}. Page 5 is cancelled by one long
% diagonal stroke and is transcribed as far as it reads; the lower half of
% p. 8 is crossed out by several strokes, likewise. Pages 11 and 12 are
% mostly diagrams (lattices of "drapeaux" D_I with arrows marked « épi »,
% small graphs); their labels and formulas are given, and each figure is
% described in a French \note, not redrawn.
%
% His pagination: none on the leaves. Numbered runs: pp. 2-4 number axioms
% « Axiome 1 », « 1' », « 2 », « 2' », « 3 », « 4 » for the « complexes
% cellulaires » (the numbering is visibly revised: an « Axiome 3 » is
% overwritten, « Axiome 4 » replaces a struck label, and p. 4 lists
% « 1, 1', 2, 2', 3, 4 »); p. 9 has conditions a)-d); p. 12 has cases
% 1)-4); p. 19 has a)-c) (c struck); p. 20 opens a numbered condition with a
% circled 1 (set as (1) with a note). No « cf. p. N » cross-reference occurs.
%
% What the batch is about. Combinatorial models of cell complexes, around
% [1977]. Pp. 2-4: « Complexes cellulaires et ensembles à opérateurs »:
% sets F_0, ..., F_n with incidence relations I_α ⊂ F_α × F_{α+1}, the
% ordered set F_* = ⊔ F_α, axioms 1-4 (projections surjective, the
% « condition du diamant »: the set of elements between two elements two
% levels apart has cardinal 2, etc.), flags (drapeaux) of type δ ⊂ [0,n],
% the set A of maximal flags and the involutions σ_α on it, the group G_n
% generated by σ_0..σ_{n-1} with σ_ασ_β = σ_βσ_α for |β-α| ≠ 1 (a Coxeter
% group of type A), the maps A/G_n^{(α)} → F_α and A/G_n^{(δ)} → A_δ, and a
% « Prop »: F_* ↦ A is fully faithful from complexes satisfying the axioms
% to G_n-sets. P. 4 ends on the geometric realisation of F_*. P. 5
% (cancelled): exact sequences of automorphism groups Aut(X, id_∂X),
% Aut(X), Aut(∂X), a covering of ∂X and Z^I. Pp. 6-8: realisation of the
% elements of F_* as cells, a Prop that for |F_*| a topological manifold
% the axioms hold (and Poincaré conjecture remarks), filtered spaces
% X_0 ⊂ ... ⊂ X_n = X, « n-variété combinatoire » as a set D_n with
% fixed-point-free involutions σ_i, consecutive ones commuting (a
% G_n-set), « n-complexe combinatoire fini » as a tower
% D_n → ... → D_0 of G_{i-1}-morphisms, its boundary δ_{i-1}, and a crossed
% out attempt at a category C_n. P. 9: « n-cartes cellulaires »: a compact
% space with a filtration by closed subspaces satisfying a)-d) (Z̄_i the
% « décomposition » of X_i along X_{i-1}, each Z̄_i^α homeomorphic to a
% ball with boundary sphere), with small figures. Pp. 11-13: lattices of
% « drapeaux » D_I for I ⊂ [0,i], épi conditions, « paracartes »
% (functors P_f^*(N) → Ens), the category Φ(C) of finite families of
% morphisms and a functor Φ(C) → Top built from cones. Pp. 15-18:
% « géométrie de drapeaux »: an ordered set D in which each D(d) =
% {d' ≤ d} is a « simplexe ordonné » P^*(I_d); equivalence with pairs (Φ,
% D : P_f^*(Φ)° → Ens) with card D({i}) = 1, strict flag geometries versus
% simplicial complexes (« ensembles simpliciaux », he writes « ens.
% quasi-simpliciaux » later), dimension functions δ : Φ → I and the case
% I = N, D[n], quasi-simplicial sets on the category of finite ordered sets
% with injective maps. P. 19: X = |D|, the closed cells |F| and open cells
% |F|°, questions on stratifications. P. 20: Δ the category of finite
% totally ordered non-empty sets with strictly increasing maps, D_* ∈ Ob Δ^
% (a presheaf), condition (1) (injectivity of Hom_Δ(Δ_n, Δ_m) → D_n),
% the ordered set D_* = ⊔ D_n, D_d ≅ P^*(Δ_m) and |D_*| = lim |Δ_m|.
%
% Notation fixed for the batch: F_* (his F with a star) for the ordered set
% of all cells; I_α for the incidences; σ_α for the involutions on maximal
% flags; G_n, G_n^{(δ)}; \mathfrak{P}_f^*(I) for his « Pf* (I) », the finite
% non-empty subsets (he writes the star or a small circle, and sometimes
% omits f); Φ for his barred-I symbol for the set of minimal elements (he
% writes it like a Φ); (Ens), Top as he writes them; D for both his roman
% and script D, which he does not seem to distinguish; \mathcal{A} for his
% script A of maximal flags.
%
% Apparatus: 99 \ill{}, 80 \uncertain{}. Hardest pages: 3 (two dense columns
% with interlinear rewrites), 6, 8 (lower half cancelled), 11-13 (diagrams).

\documentclass[11pt,a4paper]{article}
% Le préambule commun à toutes les transcriptions du fonds Grothendieck.
%
% Il tient en une idée : **l'incertitude se compose, elle ne se résout pas.**
% Une transcription qui lisse un mot douteux a détruit la seule chose pour
% laquelle on la faisait — un lecteur doit pouvoir savoir, à chaque endroit,
% ce qui était sur le feuillet et ce qui a été ajouté. Les six macros qui
% suivent sont donc l'appareil critique, et non de la décoration.
%
% Les mêmes commandes sont reconnues par `scripts/render.mjs`, qui produit la
% vue de lecture du site. Une macro ajoutée ici doit l'être aussi là-bas,
% faute de quoi le rendu échouera bruyamment — ce qui est voulu : un rendu
% silencieusement faux est pire qu'un rendu absent.

\ProvidesPackage{grothendieck}[2026/08/08 Transcriptions du fonds Grothendieck]

\usepackage{amsmath,amssymb,amsthm}
\usepackage{mathrsfs}
\usepackage{stmaryrd}
% Les diagrammes commutatifs sont partout dans ce fonds : c'est la langue
% même de la géométrie algébrique de Grothendieck, et une transcription qui
% les rendrait en prose ne transcrirait rien.
\usepackage{tikz-cd}
% Français par défaut — c'est la langue des feuillets — mais l'anglais est
% chargé aussi : la traduction et le résumé sont des documents anglais, et la
% typographie française (espaces avant les deux-points) y serait une faute.
% Ces documents basculent par \selectlanguage{english} en tête de corps.
\usepackage[english,french]{babel}
\usepackage{fontspec}
\usepackage{geometry}
\geometry{a4paper,margin=25mm,marginparwidth=28mm}
\usepackage{marginnote}
\usepackage{xcolor}
% ulem, not soul. `soul`'s \st and \ul are built on a character-by-character
% scan that XeLaTeX's font machinery breaks: the document fails with an
% undefined control sequence rather than degrading. ulem does the same two jobs
% and is Unicode-engine safe. `normalem` stops it hijacking \emph, which the
% transcriptions use for Grothendieck's own emphasis.
\usepackage[normalem]{ulem}
\usepackage{hyperref}
\hypersetup{colorlinks=true,linkcolor=black,urlcolor=black,pdfborder={0 0 0}}

% --- Métadonnées du lot -----------------------------------------------------
% Reprises telles quelles par le script de rendu, qui les affiche en tête de
% la vue de lecture. Elles fixent ce que le fichier prétend couvrir : sans
% elles, une transcription flotte sans point d'attache dans un fonds de seize
% mille feuillets.

\newcommand{\folder}[1]{\gdef\grothFolder{#1}}
\newcommand{\batch}[1]{\gdef\grothBatch{#1}}
\newcommand{\leaves}[2]{\gdef\grothFirst{#1}\gdef\grothLast{#2}}
\newcommand{\foldertitle}[1]{\gdef\grothFoldertitle{#1}}
\gdef\grothFolder{?}\gdef\grothBatch{?}\gdef\grothFirst{?}\gdef\grothLast{?}\gdef\grothFoldertitle{}

% --- Le résumé --------------------------------------------------------------
%
% Il ouvre la lecture modernisée, sous le titre « Résumé », et il est composé
% en petit. Ce n'est pas un ornement : c'est le seul passage du document qui
% ne soit pas des mathématiques, et un lecteur doit pouvoir le reconnaître
% avant de l'avoir lu — puis le sauter s'il connaît déjà le sujet, ou s'y
% arrêter s'il ne le connaît pas. Le filet qui le suit marque la reprise du
% corps.

\newenvironment{resume}
  {\small\setlength{\parskip}{0.45em}}
  {\par\medskip\hrule\medskip}

% --- Mots-clés ---------------------------------------------------------------
%
% En anglais, en clôture du résumé de la lecture modernisée : le vocabulaire
% moderne sous lequel chercher ce que les feuillets construisent. Le manifeste
% du site (`npm run manifest`) les extrait pour étiqueter la cote entière —
% cette ligne est la source unique des tags, il n'y a pas de fichier à côté.
\newcommand{\keywords}[1]{\par\smallskip\noindent
  {\footnotesize\textsc{Keywords} --- \textit{#1}}\par}

% --- L'appareil critique ----------------------------------------------------

% Le numéro de feuillet, en marge. C'est l'ancre qui permet de régler un
% désaccord sur une formule en regardant *un* feuillet plutôt qu'une chemise
% de six cents. Le site s'en sert aussi pour tourner les pages du fac-similé
% quand on fait défiler la transcription.
\newcommand{\leaf}[1]{%
  \marginnote{\footnotesize\color{gray!70}\textsf{#1}}%
  \label{leaf:#1}\ignorespaces}

% Illisible : on ne devine pas. Un mot inventé qui se lit comme les autres est
% le pire résultat possible.
\newcommand{\ill}{\textcolor{red!60!black}{[\,\ldots\,]}}

% Lecture douteuse : le mot est donné, et signalé comme incertain.
\newcommand{\uncertain}[1]{\uline{#1}}

% Ajout de l'éditeur — une lettre manquante, un mot sous-entendu.
\newcommand{\add}[1]{\textcolor{blue!50!black}{[#1]}}

% Biffé par Grothendieck lui-même. Ce qu'il a rayé fait partie du document :
% une rature dit souvent où la pensée a changé de direction.
\newcommand{\struck}[1]{\sout{#1}}

% Note du transcripteur, sur l'état matériel du feuillet.
\newcommand{\note}[1]{\par\smallskip{\small\color{gray!60!black}\textit{[#1]}}\par\smallskip}

% Note marginale de Grothendieck, distincte du corps du texte.
\newcommand{\marginal}[1]{\par\smallskip\noindent\hspace{1em}%
  {\small\color{gray!60!black}#1}\par\smallskip}

% --- Titre ------------------------------------------------------------------

\AtBeginDocument{%
  \begin{center}
    {\large\bfseries Fonds Alexandre Grothendieck --- cote n\textsuperscript{o}~\grothFolder}\\[2pt]
    {\itshape\grothFoldertitle}\\[4pt]
    {\small feuillets \grothFirst--\grothLast\ (lot \grothBatch)}\\[2pt]
    {\footnotesize\color{gray!60!black}Transcription établie d'après le fac-similé numérisé
     par l'Université de Montpellier.\\
     Les lectures incertaines sont soulignées, les illisibles notées [\,\ldots\,],
     les ajouts entre crochets.}
  \end{center}
  \vspace{1em}}

\endinput


\folder{76}
\batch{1}
\pages{1}{20}
\dating{[vers 1977]}
\watermark{Édition de démonstration}
\foldertitle{n-cartes cellulaires : notes manuscrites (s.d.)}

\begin{document}

\section*{Complexes cellulaires et ensembles à opérateurs}

\page{2}
\note{le titre est le sien, en tête de la page 2 ; il y a corrigé
« complexes » en « cartes » (surcharge au-dessus des deux occurrences), et
biffé « cellulaires » : la lecture en est \uncertain{« Cartes coniques, cartes
cellulaires, et ens. à opérateurs »}, les mots biffés restant visibles}

\struck{Complexes} \add{Cartes} \uncertain{coniques}, \struck{complexes}
\add{cartes} \struck{cellulaires}, et ens.\ à opérateurs.

$F_0,\ F_1,\ \dots,\ F_n$ \uncertain{(fini ??)}

$I_0, \dots, I_{n-1}$ ensembles ($I_\alpha$ ens.\ \ill{} $\alpha$-\ill{})
\note{la ligne commence par une suite de lettres biffée, illisible}

$\forall\, 0 \leqslant \alpha \leqslant n-1$, $I_\alpha \subset F_\alpha
\times F_{\alpha+1}$ (relation d'incidence)

\emph{Axiome 1} $I_\alpha \to F_{\alpha+1}$ \struck{\ill{}} est
\emph{surjectif}.
\marginal{\uncertain{facultatif} : Axiome 1' $I_\alpha \to F_\alpha$
surjectif}
\note{note écrite verticalement dans la marge gauche, d'une autre encre}

Sur $F_* = \coprod_{0 \leqslant \alpha \leqslant n} F_\alpha$, on met la
relation d'ordre suivante
\[
  \begin{array}{l}
  x \preceq y \ \text{ssi} \\
  x \in F_\alpha,\ y \in F_\beta
  \end{array}
  \quad
  \left\lbrace
  \begin{array}{l}
  \text{ou bien } x = y \ (\text{donc } \alpha = \beta) \\
  \text{ou bien } \alpha < \beta \text{ et } \exists\, z_\alpha = x \in
  F_\alpha,\ z_{\alpha+1} \in F_{\alpha+1}, \dots, z_{\beta-1} \in
  F_{\beta-1}, \\
  \qquad z_\beta = y \in F_\beta \text{ tels que } \forall\, \alpha
  \leqslant i < \beta,\ z_i \text{ et } z_{i+1} \text{ soient incidents}
  \end{array}
  \right.
\]

\emph{NB} Le système $(F_\alpha)_{0 \leqslant \alpha \leqslant n}$ et les
$I_\alpha$ ($0 \leqslant \alpha \leqslant n-1$) se reconstituent à partir de
l'ens.\ ordonné $F_*$, en prenant
\begin{align*}
  F_0 &= \text{ens.\ des él.\ minimaux de } F_* \\
  F_1 &= \text{\quad''\quad''\quad''\quad} F_* - F_0 \\
  &\ \ \vdots \\
  F_i &= \text{\quad''\quad''\quad''\quad} F_* - F_0 - \dots - F_{i-1}
  \qquad (i \leqslant n) \\
  I_\alpha &\subset F_\alpha \times F_{\alpha+1}, \quad I_\alpha = \lbrace
  (x,y) \mid x \in F_\alpha,\ y \in F_{\alpha+1},\ x \preceq y \rbrace
\end{align*}
[Cela montre pourquoi $I_\alpha \to F_{\alpha+1}$ surjectif]
\marginal{\emph{NB} \ill{} \ill{} \ill{}}
\note{note oblique dans la marge droite, à hauteur de ce paragraphe, en
quatre lignes serrées ; seul « NB » se lit}

Les ens.\ ordonnés ($F_*$ ainsi obtenus sont exactement \struck{\ill{}} ceux
satisfaisant la condition des chaînes, et où les chaînes (= ss-ens.\
totalement ordonnés de $F_*$) sont de \struck{les ens.\ ordonnés} \struck{les
chaînes \ill{} \ill{} de $F_*$} longueur \struck{maxima} $\leqslant n$ (on
n'exige \uncertain{pas} $F_n \neq \emptyset$).
\marginal{Vérifier}
\note{« Vérifier », à l'encre noire dans la marge gauche, avec une flèche
vers la phrase}

\struck{Ax} \emph{Drapeaux} On sait par définition les parties totalement
ordonnées non vides de $I$. Longueur d'un drapeau ($= 1 +$ cardinal), type
$\in \mathfrak{P}[0,n]$ (c'est une partie non vide de $[1,n]$).
\note{« $[0,n]$ » : le 0 est repassé à l'encre noire}
\add{\uncertain{Ens.} $D(F_*)$ des drapeaux} (\uncertain{ordonnés} par
inclusion, les él.\ de $F_*$ sont les drapeaux minimaux, la relation $x
\leqslant y$ de $F_*$ signifie \uncertain{\ill{}}
\note{la page s'arrête sur ce mot ; la phrase continue en haut de la
page 3}

\page{3}
\note{feuille à l'italienne, écrite sur deux colonnes ; on donne la colonne
de gauche, puis celle de droite}
existe un drapeau \uncertain{majorant} $x$ et $y$ \add{\ill{}} fois
\add{l'ens.\ ordonné} \struck{dans} ($F_*$ \ill{} \ill{} \ill{} l'ens.\
ordonné des drapeaux (qui a la propriété, entre autres que tout drapeau est
le sup des \ill{} \ill{} minimaux qu'il majore,
\add{et que l'ens.\ des drapeaux \ill{} \ill{} \ill{} un drapeau)}
\add{\ill{} $D(\Sigma)$}). Drapeaux maximaux… Soit $\mathcal{A}$
l'ens.\ des drapeaux maximaux.
\marginal{au \uncertain{signe} : l'ordre inverse \uncertain{puis}…}
\note{marginal à l'encre noire, oblique, en haut à gauche}

\struck{\ill{}} Pour $\forall\, x \in F_*$, soit $F(x)$ l'ens.\ des $y <
x$. \add{Il satisfait les \ill{} conditions \ill{} que $F_*$, avec $n$
\ill{} par $\alpha$ \ill{}.} \struck{On observe} l'axiome \ill{}
\note{deux lignes surchargées : une addition interlinéaire remplace
partiellement « On observe l'axiome » ; l'ordre de lecture est incertain}

\emph{Axiome} \struck{3} \emph{2} Soit $2 \leqslant \alpha \leqslant n$,
$x_\alpha \in F_\alpha$, $x_{\alpha-2} \in F(x_\alpha)_{\alpha-2}$. Alors
\struck{\ill{}} le cardinal de l'ens.\ des $x_{\alpha-1} \in
F(x_\alpha)_{\alpha-1}$ incidents à $x_{\alpha-2}$ est 2.
\marginal{\uncertain{condition} \uncertain{seulement} si $\alpha
\geqslant 2$}

\emph{Axiome 3} Soit $x_1 \in F_1$, alors le cardinal de l'ens.\ des
$y_0 \in$ \struck{\ill{}} $F(x)_0$ \struck{incidents à $x$} est 2.
\marginal{\ill{} $\alpha \geqslant 1$}

\struck{Définition}

Soit donc $\mathcal{A}$ l'ens.\ des drapeaux maximaux. \struck{\ill{}} On
va introduire sur $\mathcal{A}$ des opérations $\sigma_\alpha$ ($0
\leqslant \alpha \leqslant n-1$). Soit $d = (x_0, x_1, \dots, x_\nu)$ ($\nu
= n$) un drapeau maximal, avec $x_i \in F_i$ ($0 \leqslant i \leqslant
\nu$).
\[
  \sigma_0(d) = (x_0', x_1, \dots, x_\nu)
\]
\struck{où $x_0'$ est} caractérisé par la condition $x_0' \neq x_0$ [et
$(x_0', x_1, \dots, x_\nu)$ un drapeau, i.e.\ $x_0' \in F(x_1)_0$] cf
Axiome 2
\[
  \sigma_\alpha(d) = (x_0, \dots, x_{\alpha-1}, x_\alpha', x_{\alpha+1},
  \dots)
\]
caractérisé par $x_\alpha' \neq x_\alpha$ (\emph{et} par
$\sigma_\alpha(d)$ est un drapeau)
\marginal{$1 \leqslant \alpha \leqslant n-1$}
\marginal{\uncertain{le rendre naturel}}
\note{la note « le rendre \uncertain{naturel} » est écrite en oblique entre
les deux colonnes}

Considérons aussi l'axiome (facultatif)

\emph{Axiome 2'} $\forall\, x_{n-1} \in F_{n-1}$, le cardinal de l'ens.\
des \struck{drapeaux} $y \in F_n$ incidents à $x_{n-1}$, est 2.
On pose alors
\[
  \sigma_n(d) = (x_0, \dots, x_{n-1}, x_n') \quad \text{caractérisé par }
  x_n' \neq x_n \quad (\text{etc.}\dots)
\]
\note{après $(x_0, \dots, x_{n-1}, x_n')$, quelques mots biffés, dont
« $x_n'$ est le » ; « caractérisé par $x_n' \neq x_n$ (etc…) » est
écrit au-dessus}
\marginal{\uncertain{variante} \ill{}}

\emph{NB} Les $F(x)$ (où $x \in F_\alpha$) satisfont les axiomes 1, 2, 3,
\emph{et 4} avec $n$ remplacé par $\alpha - 1$…

\note{colonne de droite}
On va voir maintenant que (moyennant les axiomes $1, 1', 2, 2', 3$
\add{et un axiome 4 plus bas}) l'ens.\ à opérateurs $\mathcal{A}$ permet
de \emph{reconstituer} $F_*$ (donc $(F_\alpha)$, $(I_\alpha)$…).

\struck{On a évidemment} \struck{Les} Mais donnons d'abord les relations de
commutation évidentes
\[
  \sigma_\alpha \sigma_\beta = \sigma_\beta \sigma_\alpha \quad \text{si }
  |\beta - \alpha| \neq 1 \ \text{i.e.\ } \alpha, \beta \text{ non
  consécutifs.}
\]
Soit $G_n$ le \emph{groupe engendré} par les \struck{relations} gén.\
$\sigma_\alpha$ ($0 \leqslant \alpha \leqslant n-1$) avec les relations
précédentes. $\mathcal{A}$ est un $G_n$-ensemble. Les $\sigma_\alpha$
($0 \leqslant \alpha \leqslant n-1$) opèrent sans pt fixe. Soit
$G_n^{(\alpha)}$ le sous-groupe de $G_n$ engendré par les $\sigma_\beta$
avec $\beta \neq \alpha$. On a \uncertain{évidemment}
\[
  \mathcal{A}/G_n^{(\alpha)} \longrightarrow F_\alpha \qquad (0 \leqslant
  \alpha \leqslant n-1)
\]
\emph{Axiome 4} \note{« Axiome 4 » est écrit au-dessus d'un mot
lourdement biffé} Cette application est \emph{bijective}.

(La surjectivité est claire \add{par l'axiome 1}. Pour la bijectivité, il
suffit de voir que si $d = (x_0, \dots, x_n)$, $d' = (x_0', \dots, x_n')$
sont deux $\in \mathcal{A}$ avec \struck{\ill{}} $x_\alpha = x_\alpha'$,
alors on peut passer de $d$ à $d'$ en appliquant successivement des
transformations $\sigma_\beta$ avec $\beta \neq \alpha$.

Plus généralement, soit $\delta \subset [0,n]$, \add{$\delta \neq
\emptyset$}, et considérons le sous-groupe $G_n^{(\delta)}$ des $G_n$
engendré par les $\sigma_i$ avec $i \in \complement_{[0,n]} \delta$.
Alors on a
\[
  \mathcal{A}/G_n^{(\delta)} \longrightarrow \mathcal{A}_\delta \quad
  (\text{drapeaux de type } \delta)
\]
et on voudra que ce soit bijectif (\uncertain{conséquence} \ill{} de
l'axiome 4 et des précédents?))

\page{4}
Ceci dit, les relations d'inclusion naturelles entre les $G_n^\delta$
($G_n^\delta \subset G_n^{\delta'}$ pour $\delta' \subset \delta$)
définissent les relations d'incidence \add{(= \uncertain{relations
d'inclusion}) entre} drapeaux de types divers. On reconstitue ainsi
l'ens.\ ordonné des drapeaux des divers types $\delta$, et par là l'ens.\
des $F_\alpha$ (drapeaux minimaux) et leurs relations d'incidence…

De façon précise :

\emph{Prop} Le foncteur $(F_*) \mapsto (\mathcal{A})$ des complexes
cellulaires satisfaisant les axiomes $1, 1', 2, 2', 3, 4$
(\uncertain{confirmé ?}) vers les ens.\ à groupe d'opérateurs $G_n$, est
\emph{pl.\ fidèle}.

Il faudrait déterminer l'image essentielle.

\note{un trait de séparation horizontal, barré d'un trait vertical}

Réalisation géom.\ d'un complexe $F_*$ (satisfaisant le seul axiome 1) :
c'est par définition la réalisation géométrique de la triangulation
simpliciale définie par l'ens.\ des parties \add{(finies)} totalement
ordonnées de $F_*$.

\page{5}
\note{toute la page est barrée d'un long trait diagonal ; elle se lit
néanmoins. Ce sont des suites exactes empilées, sans texte ; on en donne
les lignes, flèches verticales comprises dans une note}
\[
  1 \to \underline{\mathrm{Aut}}(X, \mathrm{id}_{\partial X}) \to
  \mathrm{Aut}(X) \to \mathrm{Aut}(\partial X) \to 1
\]
\note{au-dessus de $\mathrm{Aut}(X)$ une flèche monte vers $1$ ;
au-dessous, une colonne $1 \to \mathbf{Z}^I \to
\widetilde{\mathrm{Aut}}(X) \to \mathrm{Aut}(X)$, une flèche horizontale
arrivant de gauche sur $\widetilde{\mathrm{Aut}}(X)$}

\begin{tikzcd}
  \widetilde{\partial X} \arrow[d, "p"] & \\
  \partial X \arrow[r, hook'] & X
\end{tikzcd}

\[
  \widetilde{\mathrm{Aut}}(X) = \lbrace (u, \tilde{v}) \mid u \in
  \mathrm{Aut}(X),\ \tilde{v} \in \mathrm{Aut}(\widetilde{\partial X}),\ \ 
  (u \mid \partial X) \circ p = p \circ \tilde{v} \rbrace \subset
  \mathrm{Aut}(X) \times \mathrm{Aut}(\widetilde{\partial X})
\]
\note{la flèche d'inclusion $X \hookleftarrow \partial X$ est tracée de
droite à gauche}

\[
  \begin{array}{ccccccccc}
    1 & \to & S\mathcal{A}(X) & \to & \mathcal{A}(X) & \to &
    \mathcal{A}(\partial X) & \to & 1 \\
    & & \uparrow & & \uparrow & & \uparrow & & \\
    1 & \to & S\widetilde{\mathcal{A}}(X) & \to &
    \widetilde{\mathcal{A}}(X) & \to & \mathcal{A}(\widetilde{\partial X})
    & \to & 1 \\
    & & \uparrow & & \uparrow & & \uparrow & & \\
    1 & \to & \mathbf{Z}^I & \Rightarrow & \mathbf{Z}^I & \to & 1 & \to &
    1
  \end{array}
  \qquad I = \pi_0(\partial X)
\]
\note{en tête des deux premières colonnes, des flèches pointillées venues
de $1$ ; une flèche pointillée va en diagonale de $S\mathcal{A}(X)$ vers
$\widetilde{\mathcal{A}}(X)$ ; sous $S\widetilde{\mathcal{A}}(X)$ il
écrit « $\simeq S\mathcal{A}(X) \times \mathbf{Z}^I$ », sous la troisième
colonne une flèche descend de $1$. La flèche « $\Rightarrow$ » est son
signe d'égalité fléchée}

\[
  \begin{array}{ccccccccc}
    1 & \to & ST(X) & \to & \mathcal{A}(X)/S\mathcal{A}^\circ(X) & \to &
    \mathcal{A}(\partial X) & \to & 1 \\
    & & \uparrow & & \uparrow & & \uparrow\!\wr & & \\
    1 & \to & \widetilde{ST}(X) & \to & \widetilde{\mathcal{A}}(X)/
    S\mathcal{A}^\circ(X) & \to & \mathcal{A}(\widetilde{\partial X}) &
    \to & 1 \\
    & & \uparrow & & \uparrow & & \uparrow & & \\
    1 & \to & \mathbf{Z}^I & \Rightarrow & \mathbf{Z}^I & \to & 1 & \to &
    1
  \end{array}
\]
\note{même disposition, avec la flèche pointillée de $ST(X)$ vers
$\widetilde{\mathcal{A}}(X)/S\mathcal{A}^\circ(X)$, et sous
$\widetilde{ST}(X)$ l'annotation « $\simeq ST(X) \times \mathbf{Z}^I$ » ;
le dénominateur $S\mathcal{A}^\circ(X)$ est \uncertain{lu} ainsi dans les
deux lignes}

\page{6}
Les éléments de $F_*$ se réalisent comme \uncertain{des} parties
\add{(i.e.\ structures coniques)} de $|F_*|$ [\uncertain{associé} à $x
\in F_\alpha$ correspond la réunion des simplexes réalisation des simpl.\
\struck{formés} \add{ouverts} d'éléments $\leqslant x$], la relation
d'incidence devient l'inclusion --- il y a une équivalence de catégories
(\struck{\ill{}} d'objets top.\ versus objets combinatoires) à \uncertain{dégager}.

\emph{Prop} Si $|F_*|$ est une variété top.\ \add{\uncertain{pointée}}
\uncertain{de dimension $n$}, alors $(F_*)$ satisfait les conditions
$1, 1', 2, 2', 3, 4$… [Plus précisément, pour que $(F_*)$ soit une
variété top.\ de dim.\ $n$, il faut et il suffit que $(F_*)$ satisfasse
$1, 1', 2, 2', 3, 3'$, $4$, \emph{et} que $\forall\, x \in F_\alpha$
\add{($\alpha \geqslant 3$)}, la réalisation géom.\ de $F(x)$ soit une
sphère [en fait, \struck{\ill{}} \uncertain{modulo} l'hyp.\ de Poincaré,
il suffit d'exiger que \struck{pour $\alpha \geqslant 3$}, $|F(x)|$ soit
simpl.\ connexe et \add{pour $\alpha \geqslant 5$, elle ait de plus}
\struck{\ill{}} l'homologie d'une sphère de même dimension.
\note{« $\alpha \geqslant 3$ » est corrigé en « $\geqslant 5$ » dans
l'addition ; le détail des surcharges n'est pas sûr}

\note{trait de séparation}

Soit $X$ un espace top.\ (filtré par \add{(compact ??)}
\[
  X_0 \subset X_1 \subset \dots \subset X_n = X .
\]
\struck{On suppose} \ill{} que \uncertain{$\forall\, i$}, la
\emph{différence} de $X_i$, \uncertain{moins} $X_{i-1}$, soit
\struck{homéomorphe à une} \ill{} \struck{bord de} \ill{} \ill{} de boules
$B^{i}$. Soit $\mathcal{A}$ l'ens.\ des « drapeaux » relatifs à cette
filtration (à définir avec soin). On a alors une op.\ de $G_n$
\add{(\struck{\ill{}} $= G_n$ si $X$ \uncertain{variété} top.\ de
\uncertain{dim.\ $n$})}.
\note{la fin de la page est très serrée et partiellement surchargée}

\page{7}
avec les $\sigma_i$ opérant sans pt fixe.

\emph{Prop} Le foncteur de la catégorie \struck{\ill{}} des \add{variétés
de dim.\ $n$} $X$ filtrées comme dessus, vers les $G_n$-ensembles finis
(où les $\sigma_i$ opèrent sans pt fixe) est pl.\ \emph{fidèle}.

Déterminer l'image essentielle (\uncertain{moyennant} l'hyp.\ de Poincaré
au besoin) et montrer qu'elle est plus grande que celle correspondant aux
complexes cellulaires se réalisant suivant une variété topologique…

\emph{$n$-Variété combinatoire} \add{finie} définie par un ens.\
\add{fini} $D_n$ muni de $n+1$ \uncertain{automorph.} involut.\ sans pt
fixe $\sigma_i$ ($0 \leqslant i \leqslant n$) tels que deux $\sigma_i$ non
consécutifs commutent --- i.e.\ $D_n$ est un ens.\ à op.\ sous $G_n$, avec
les $\sigma_i$ ($0 \leqslant i \leqslant n$) op.\ sans pt fixe.

\emph{$n$-complexe combinatoire} \emph{fini}. Donné par un système
$D_n, D_{n-1}, \dots, D_1, D_0$ d'ens.\ \add{finis}, \struck{chaque}
\add{$D_i$ ($1 \leqslant i \leqslant n$)} étant un \struck{\ill{}}
$G_{i-1}$-ensemble \add{« spécial »} (i.e.\ les gén.\ $\sigma_\alpha$ op.\
sans pt fixe) et des applications
\[
  D_n \longrightarrow D_{n-1} \longrightarrow D_{n-2} \longrightarrow
  \cdots \longrightarrow D_1 \longrightarrow D_0
\]
telles que $\forall\, 2 \leqslant i \leqslant n$, $D_i \to D_{i-1}$ soit
un $G_{i-2}$-morphisme, et que pour \struck{$0 \leqslant j \leqslant$}
$i \geqslant j+2 \geqslant j \geqslant 0$,

\page{8}
\note{feuille à l'italienne ; seule la moitié gauche est écrite}
l'application composée $D_i \to D_j$ (qui est un $G_{j-1}$-morphisme) soit
invariante par $\sigma_{i-1}, \dots, \sigma_{j+1}$.

\emph{$n$-variété combinatoire} $\Leftrightarrow$ $G_n$-ens.\ $D_n$. Si
$0 \leqslant i \leqslant n$, le \add{bord de la $i$-décomp.}
\struck{\ill{} décomp.} $\partial(\mathrm{D\acute{e}c}(X_i/X_{i-1}))$ est
une $(i-1)$ variété combinatoire $\delta_{i-1}$ qui est
\begin{align*}
  \delta_{n-1} &= (D_n;\ \sigma_0, \dots, \sigma_{n-1}) \\
  \delta_{n-2} &= (D_n/\sigma_n;\ \sigma_0, \dots, \sigma_{n-2}) \\
  \delta_{n-3} &= (D_n/(\sigma_n, \sigma_{n-1});\ \sigma_0, \dots,
  \sigma_{n-3}) \\
  &\ \ \cdots \\
  \delta_{i-1} &= (D_n/(\sigma_n, \dots, \sigma_{i+1});\ \sigma_0, \dots,
  \sigma_{i-1})
\end{align*}

\note{la suite de la page est barrée de plusieurs grands traits obliques ;
on la donne telle qu'elle se lit}
\struck{On définit par récurrence} \struck{sur $n$ la cat.\
$\mathcal{C}_n$ des} $n$-\add{cartes} \struck{\ill{}} \struck{combinatoires}
\add{\uncertain{finies} et foncteur $i_n : \mathcal{C}_n \to
\mathcal{V}_n$}, \struck{en posant}
\note{sous $\mathcal{V}_n$, en petit : « $n$-var.\ comb. »}

\struck{$\mathcal{C}_0 =$ \add{$0$-sphères} \ill{} $=$ ens.\ finis}

\struck{$\mathcal{C}_n = \lbrace X_{n-1}, S_{n-1}, \varphi \rbrace$}
\add{\struck{cat.\ des}}

\struck{où $X_{n-1} \in \mathrm{Ob}\,\mathcal{C}_{n-1}$,} \struck{$S_{n-1}
\in \mathrm{Ob}\,\mathcal{V}_{n-1}$,} \struck{$\varphi : i_{n-1}(S_{n-1})
\to X_{n-1}$}

\struck{$i_n : \mathcal{V}_n \to \mathcal{C}_n$ défini par $i_n(D_n) =
($} \note{la phrase s'arrête sur la parenthèse ; en dessous, « $G_n$-ens. »
; la flèche $i_n$ est écrite ici dans l'autre sens que plus haut}

\page{9}
\struck{Complexes} \add{Cartes} coniques, \struck{complexes} \add{cartes}
\struck{cellulaires}, et ens.\ à opérateurs.

Une \emph{$n$-carte conique} \add{(finie)} est un espace top.\ $X$
compact, muni d'une filtration croissante \add{par des parties fermées}
\[
  X_0 \subset X_1 \subset \dots \subset X_{n-1} \subset X_n = X
\]
telles que
\begin{itemize}
  \item[a)] $X_0$ fini (i.e.\ discret)
  \item[b)] pour $1 \leqslant i \leqslant n$, \add{soit}
  $\mathrm{D\acute{e}c}(X_i/X_{i-1}) = \overline{Z}_i$ (découpage de $X_i$
  suivant $X_{i-1}$), $\dot{Z}_i = \overline{Z}_i \mid X_{i-1}$
  (\emph{NB} on a $\overline{Z}_i \to X_i$, induisant un isom.\
  $\overline{Z}_i - \dot{Z}_i \to X_i - X_{i-1}$), alors
  \struck{$\overline{Z}_i$ a un \ill{}} \add{\ill{}} fini de comp.\
  connexes, $\overline{Z}_i^\alpha$, et chacune est homéom.\ au cône sur
  $\dot{Z}_i^\alpha = \overline{Z}_i^\alpha \cap \dot{Z}_i$ (et
  $\dot{Z}_i^\alpha \neq \emptyset$)
  \item[c)] $\dot{Z}_i \to X_{i-1}$ est un morphisme fini, \struck{dont
  l'image} [$\Leftrightarrow \dot{Z}_i \to X_{i-1}$ \ill{} \ill{}
  \ill{}]
  \item[d)] La filtration de $X$ est équisingulière \ill{} \ill{}
\end{itemize}
\note{d) est récrit à droite en surcharge : « d) Soit $X_{i-1}^!$ l'image
de $\dot{Z}_i$ dans $X_{i-1}$ ; \struck{\ill{}} Alors $X_{i-2}^! =
X_{i-1}^! \cap X_{i-2}$ » (lecture \uncertain{incertaine} des indices)}

Cas $n = 0$ : espaces finis discrets.

Cas $n = 1$ : $X = X_1 \supset X_0$. \add{\uncertain{cpct}}
\[
  \dot{Z}_1 = \coprod_\alpha \dot{Z}_1^\alpha \to X_0
\]

\note{au bas de la page, à gauche, six petites figures : une boucle
posée sur un point, un segment, un segment à deux sommets, un
« Y » à trois sommets cerclés, un segment isolé, un disque bordé d'une
boucle avec un point, et un sommet d'où partent trois arêtes, l'une
hachurée ; ce sont des exemples de 1-cartes. Sous elles,
« $\dot{Z}_1 \subset \overline{Z}_1$ » au-dessus de « $X_0 \subset X_1$ »,
avec deux flèches verticales}

\note{à droite, des tableaux d'inclusions : « $x_0 \in x_1$, $x_1$,
$x_2$ » au-dessus de
$\dot{Z}_1 \subset \overline{Z}_1$, $\dot{Z}_2 \subset \overline{Z}_2$,
$\dot{Z}_3 \subset \overline{Z}_3$, qui s'envoient sur $X_0 \subset X_1
\subset X_2 \subset X_3$ ; en dessous $x_0 \subset x_1 \subset x_2 \subset
x_3$ ; puis le schéma général ci-dessous}
\[
  \begin{array}{ccccccc}
    x_{i+1} & & \dot{Z}_{i+1} \subset Z_{i+1} & & \dot{Z}_{i+2} \subset
    \overline{Z}_{i+2} & & \\
    & & \downarrow & & \downarrow & & \\
    X_{i-1} & \subset & X_i \subset X_{i+1} & & \subset & & X_{i+2}
  \end{array}
\]
\[
  \begin{array}{ccccc}
    Z_{i+1} \mid X_{i-1} & \subset & Z_{i+1} \mid X_i & \to &
    \overline{Z}_{i+2} \\
    \downarrow & & \downarrow & & \downarrow \\
    X_{i-1} & \subset & X_i & & X_{i+2}
  \end{array}
\]
\note{les deux derniers tableaux sont serrés et en partie surchargés ;
les indices sont \uncertain{lus} au mieux}

\page{11}
Ens.\ ordonné des cellules \struck{\ill{}} ass.\ $X_i$

\[
  \begin{array}{c}
    \overline{X}_{i+1} \\
    \mid \\
    X_{i-1}\ X_i
  \end{array}
  \qquad X_i
\]

Géométrie des drapeaux de dim $\leqslant n$

$D_0$ \quad \emph{Ex} $n = 0$ \quad $D_0$

\quad\ $n = 1$ \quad $D_{01}$ \emph{épi !!} $\qquad D_0 \leftarrow D_{01}
\to D_1$

\quad\ $n = 2$
\note{à gauche, deux essais biffés avec $X_{i+1}$, $X_{i-1}$, $X_i$}

\note{figure : le treillis des drapeaux pour $n = 2$. En haut $D_{012}$ ;
de lui descendent des flèches vers $D_{01}$, $D_{02}$, $D_{12}$, marquées
« épi » ; de ceux-ci vers $D_0$, $D_1$, $D_2$, la flèche $D_{02} \to D_1$
barrée d'une croix, et « épi » sur plusieurs flèches}

\[
  D_{01\dots i} \to D_I \ \text{\emph{épi}} \quad \text{si } i = \sup I
  \quad \text{i.e.\ } i \in I \subset [0,i]
\]

$D_I$

\note{figure centrale : le même treillis dessiné plus grand, $D_{012}(F)$
en haut, $D_{02}(F)$, $D_{12}(F)$, puis $D_{01}$, $D_0$, $D_1$ et une
case $\lbrace 1 \rbrace$ ; au-dessous une petite échelle encadrée « 2 3 4
5 ». À droite, un carré $D_{012} \to D_{12}$, $D_{02} \to D_2$, puis deux
segments dessinés, l'un avec un point milieu, et « 23 ? »}

$D_\sigma \to D_{\sigma'}$ \emph{épi} si $\sigma' \subset \sigma$ et
… $\sigma$ \ill{}, i.e.\ $\sigma = \sigma'' \cup \sigma'$, $\sigma''
< \sigma'$

\note{à gauche, « $\exists$ $d'$ » avec $D_{123} \to D_{13}$ et
$D_{123} \to D_{13}$ (second indice \uncertain{lu} $123$) ; au-dessous, deux treillis pour $n = 3$
: le premier, $D_{012}$ au sommet, $D_{01}$, $D_{02}$, $D_{12}$, $D_0$,
$D_1$, $D_2$, avec des flèches épaisses et « épi » ; le second, $D_{0123}$
au sommet, $D_{013}$, $D_{023}$, $D_{123}$, $D_{03}$, $D_{13}$,
$D_{23}$, $D_3$, les flèches vers $D_3$ repassées en gras, et un
troisième treillis en pointillé relié au premier par de longues
flèches}

\begin{align*}
  D_{01} &\to D_1 & D_{013} &\to D_{13} \\
  D_{012} &\to D_{12} & D_{0123} &\to D_{123} \\
  D_{12} &\to D_2 & D_{123} &\to D_{23} \\
  D_{02} &\to D_2 & D_{23} &\to D_3
\end{align*}

\note{en bas à droite, une figure : un arc hachuré de lignes parallèles
numérotées $0, 1, 2, 3$, avec $D_1$, $D_2$ et une flèche vers $D_3$ ;
au-dessous « $1$, $1 + 1$, $1 + 2$ » avec « $\shortparallel\ 2$ »}
\[
  n_i = \sum_{j < i} n_j + 1
\]

\page{12}
\note{en haut à gauche, quatre cas numérotés}
1) $F_0$ \qquad 2) $F_0 \leftarrow R_1 \to F_1$, épi \qquad 3) \struck{un
premier essai hachuré} avec $F_0$, $F_1$, $F_2$

\note{figure : un diagramme à deux étages. En haut $R_1'$, avec des
flèches vers $F_0'$ et $F_1'$ ; de $F_0'$, $R_1'$, $F_1'$ descendent des
flèches vers $R_1$, $F_0$, $F_1$, $F_2$ ; $R_1$ s'envoie sur $F_0$ et
$F_1$. Le même diagramme est répété plus bas pour le cas 4), précédé de
« $F_0$ »}

\note{en haut à droite}
\emph{\uncertain{décomposition} cellulaire} \struck{\ill{}} \add{\ill{}
\ill{}} : l'ens.\ \struck{\uncertain{quasi-simplicial}} \struck{\ill{}}
\struck{$F_1 = F_0 \amalg F_1$}, ordonné des drapeaux
\begin{align*}
  \text{type } 0 &: F_0 \\
  \text{type } 1 &: F_1 \\
  \text{type } (0,1) &: R_1
\end{align*}
\note{figure : un graphe, deux sommets cerclés reliés par une longue
arête ; au premier pendent trois arêtes courtes ; le second, cerclé en
bas, porte deux arêtes marquées d'un point milieu et une arête vers un
troisième sommet cerclé, d'où partent trois arêtes formant un fuseau}

$\Phi(\mathcal{C}) =$ Catégorie des \struck{\ill{}} \add{familles
finies} $(f_i)_{i \in I}$ de morphismes de $\mathcal{C}$ de même but
\note{à gauche, biffé : « $\mathcal{C}_0 = \mathrm{Ens}_f$ » et
« $\mathcal{C}_1 =$ »}

Un foncteur $\mathcal{C} \xrightarrow{\ \varphi\ } (\mathrm{Top}$ \struck{\ill{}}$)$ implique un foncteur
\[
  \Phi(\mathcal{C}) \xrightarrow{\ \Phi(\varphi)\ } \mathrm{Top}
\]
\[
  (X_i \xrightarrow{f_i} X)_{i \in I} \longmapsto
  \coprod_i \mathrm{C\hat{o}ne}\, \varphi(X_i)
  \ \amalg_{\coprod \varphi(X_i)}\ \varphi(X)
\]
\note{la lettre avant $(X_i)$ est surchargée, lue $\varphi$ ; l'indice
sous le second $\amalg$ est $\coprod \varphi(X_i)$}

Si les $\varphi(f)$ ($f \in \mathrm{Fl}(\mathcal{C})$) sont finis, de
même les $\Phi\varphi(f)$, $f \in \mathrm{Fl}(\Phi\mathcal{C})$ ; si
$\varphi$ est \struck{\ill{}} fidèle \uncertain{(resp.\ pl.\ fidèle)}
$\Phi\varphi$, itou vers les espaces top.\ filtrés, \emph{mod isotopie}…
\[
  \Phi(\Phi(\mathcal{C})) \to \mathrm{Top}
\]

\page{13}
\struck{\ill{}} \emph{paratriangulation}

\emph{paracartes} : foncteurs $\mathfrak{P}_f^*(\mathbf{N})
\xrightarrow{D} (\mathrm{Ens})$ tels que \uncertain{$\sigma \subset
\sigma'$} \add{segment finis} $\Rightarrow D(\sigma) \to D(\sigma')$
épi.

\emph{paracartes uniquedimensionnelles} : $D(\sigma) \to D(\sigma')$ bij.\
\uncertain{épi} si \struck{\ill{}} $\sigma \subset \sigma' \subset
\Delta_n$, et $D(\lbrace i \rbrace) = 0$ si $i > n$.

\emph{paracartes} : \struck{\ill{}} \emph{stratifications
équidimensionnelles}\note{le premier $\sigma$ de la formule est surchargé,
peut-être sur un $\Delta$ ; lecture \uncertain{incertaine}} : $D(\sigma) \to D(\sigma')$ épi si $\sup \sigma = \sup \sigma'$, i.e.\ $D([0,i]) \to
D(\sigma)$ épi si $i = \sup \sigma$.

\emph{paracartes spéciales} :
\uncertain{les cellules} sont équi-dim ; $\forall$ \struck{\ill{}}
$D(0, \dots, j-1, j+1, \dots, i) \leftarrow D(0, \dots, i)$ de degré 2.
\note{ces deux conditions sont réunies par une accolade}
\note{la seconde ligne de l'accolade porte, biffé, un premier essai
contenant « $D(\lbrace i \rbrace) \neq \emptyset$ » ; « de degré 2 » est
souligné}

\note{à gauche, un petit diagramme : $D_{01} \to D_1$ et $D_{012}$ avec
deux flèches vers $D_{12}$ et $D_{02}$}

\emph{Question} : \ill{} \ill{} $D(\dots, j-1, j, j+1, \dots) \to
D(\dots, j-1, j+1, \dots)$ \uncertain{toujours} de degré 2 ??

p.\ ex.\ $D_{123} \to D_{13}$ degré 2…

d'où groupe $G_i$ \uncertain{engendré} par $\sigma_0, \dots, \sigma_{i-1}$
\uncertain{opérant sur} $R_i \shortparallel D(0, \dots, i)$.

On veut conditions \uncertain{assurant} que $D(0, \dots, i)$
\emph{\uncertain{récupère}} tous les $D(i_0, \dots, i_p)$ \uncertain{via}
$(R_i, G_{i-1})$ et les applications $R_{i+1} \to R_i \to \cdots$

\page{15}
Un ens.\ ordonné $D$ est appelé « \emph{géométrie de drapeaux} » si pour
tout $d \in D$, l'ens.\ ordonné \add{\emph{induit}} $D(d) = \lbrace d' \in
D \mid d' \leqslant d \rbrace$ a la structure d'un « simplexe ordonné »
i.e.\ est isomorphe à $\mathfrak{P}^*(I_d)$,
\marginal{parties finies $\neq \emptyset$}
\marginal{\emph{NB} on \uncertain{demande} $I_d \neq \emptyset$}
avec l'inclusion --- où $I_d$ peut s'identifier à l'ens.\ des él.\
minimaux de $D(d)$, on encore : la partie de l'ens.\ $\Phi = D(1)$
\note{« $D(1)$ » : le 1 est cerclé, peut-être un indice de degré} des
él.\ minimaux de $D$ (« drapeaux de longueur 1 » --- \emph{NB}
$\mathrm{card}(I_d)$ s'appellera « la longueur » du drapeau $d$). Pour
toute partie finie $\sigma$ de $\Phi$, soit $D(\sigma)$ l'ens.\ des $d \in
D$ tels que \struck{\ill{}} $I_d = \sigma$, on trouve un
contrafoncteur
\[
  (\text{parties finies} \neq \emptyset \text{ de } \Phi)^\circ
  \xrightarrow{\ D\ } (\mathrm{Ens})
  \qquad \text{tel que } \mathrm{card}\, D(\lbrace i \rbrace) = 1 \quad
  \forall\, i \in \Phi
\]
dont la connaissance permet de reconstituer $D$ comme
\[
  D \simeq \coprod_{\sigma \in \mathfrak{P}_f^*(\Phi)} D(\sigma)
\]
avec la relation d'ordre évidente. Ainsi
\[
  \left\lbrace
  \begin{array}{l}
    \text{géométries de drapeaux} \\
    \text{(ens.\ ordonnés ainsi)}
  \end{array}
  \right.
  \Longleftrightarrow
  \begin{array}{l}
    \text{Couples } (\Phi, D) \text{ d'un ens.\ } \Phi \text{ et d'un
    foncteur} \\
    \mathfrak{P}_f^*(\Phi)^\circ \longrightarrow (\mathrm{Ens}) \\
    \text{tels que } \mathrm{card}\, D(\lbrace i \rbrace) = 1\ \forall\, i
    \in \Phi
  \end{array}
\]
\note{« ainsi » dans la parenthèse de gauche est une lecture
\uncertain{incertaine}}
\marginal{\emph{NB} $\Phi$ est aussi l'ens.\ des sommets d'un
\uncertain{ens.\ quasi-simplicial}, mais \ill{} \ill{} \ill{} $D$ n'est pas
déterminé par celui-ci…}

\emph{NB} Pour que dans $D$ on ait $\forall\, d \in D$
\[
  d = \sup_{i \in I_d} i
\]
\quad $\Big|$ géométrie de drapeaux stricte

il faut et il suffit que $D$ soit isomorphe à un ens.\ de parties finies
\add{$\neq \emptyset$} de $\Phi$ avec l'inclusion (contenant, pour toute
partie finie \uncertain{$\sigma$} qu'elles contient, ses sous-ens.\ $\neq
\emptyset$ plus petits). On trouve ainsi :
\[
  \left(
  \begin{array}{c}
    \text{Géom.\ de drapeaux} \\
    \text{strictes}
  \end{array}
  \right)
  \Longleftrightarrow
  \left(
  \begin{array}{c}
    \text{ensembles} \\
    \text{simpliciaux}
  \end{array}
  \right)
\]
\note{« ensembles » est une lecture \uncertain{incertaine} du premier mot
de la parenthèse de droite}
En termes du foncteur $D : \mathfrak{P}_f^*(\Phi)^\circ \to
(\mathrm{Ens})$, cela signifie que $D$ est \struck{\ill{}} \add{un
sous-foncteur du} foncteur \struck{\ill{}} \add{final} \struck{\ill{}}
i.e.\ $\forall\, \sigma$, \struck{\ill{}} $\mathrm{card}\, D(\sigma)
\leqslant 1$ i.e.\ $\mathrm{card}\, D(\sigma) \in \lbrace 0, 1 \rbrace$.
\note{fin de page très surchargée : les additions interlinéaires
« un sous-foncteur du » et « final » sont \uncertain{lues}}

\page{16}
\add{Soit $I$ un ens.\ \ill{}}
Supposons qu'on ait un foncteur
\[
  \mathfrak{P}_f^*(I)^\circ \xrightarrow{\ D'\ } (\mathrm{Ens})
\]
sans condition de la forme $\mathrm{card}\, D(\lbrace i \rbrace) = 1 \quad
\forall\, i \in I$. On peut néanmoins lui associer une géométrie de
drapeaux \struck{dont l'ens.\ des}
\[
  D = \coprod_{\sigma \in \mathfrak{P}_f^*(I)} D'(\sigma)
\]
dont l'ens.\ des él.\ minimaux est
\[
  \Phi = \coprod_{i \in I} D'(\lbrace i \rbrace)
\]
\note{sous ce signe somme, un indice biffé, remplacé par « i ∈ I »}
Il a donc une application
\[
  \Phi \xrightarrow{\ \delta\ } I
\]
(application « dimension » $\delta$ dans les cas qui nous intéressent)
ayant la propriété que pour tout $d \in D$, la restriction
\note{ici un premier essai de la formule, biffé et en partie illisible}
\[
  \delta \mid I_d : I_d \to I
\]
est injective. Inversement, si on a une géométrie de drapeaux $D$ et une
application $\delta$ de $\Phi = D(1)$ (ens.\ des él.\ minimaux de $D$)
dans un ens.\ $I$, satisfaisant la condition précédente, on lui associe un
foncteur
\[
  D' : \mathfrak{P}_f^*(I) \to (\mathrm{Ens}) \quad \text{par}
\]
\[
  D'(\sigma) = \coprod_{\substack{\sigma' \in \mathfrak{P}_f^*(\Phi) \\
  \text{t.q.\ } \delta \mid \sigma' : \sigma' \xrightarrow{\sim} \sigma}}
  D(\sigma')
\]
et on voit que $(D, I, \delta)$ se reconstitue à partir de $(I, D')$.
\ill{}

\page{17}
\[
  \left.
  \begin{array}{l}
    \text{Géom.\ de drapeaux } D \\
    + \text{« fonction dimension »} \\
    \delta : \Phi \longrightarrow I \\
    \quad \text{\scriptsize él.\ min.}
  \end{array}
  \right|
  \Longleftrightarrow
  \left|
  \begin{array}{l}
    \text{foncteur } D' : \mathfrak{P}_f^*(I)^\circ \to (\mathrm{Ens}) \\
    \text{\emph{NB} Posant } I_0 = \lbrace i \in I \mid D'(\lbrace i
    \rbrace) \neq \emptyset \rbrace \\
    \text{on a sur } I_0 \text{ une structure } \dots
  \end{array}
  \right.
\]
\note{dans l'encadré de gauche, « fonction » (lecture
\uncertain{incertaine}) est écrit au-dessus d'un mot biffé illisible,
\struck{\ill{}}}
\note{la fin de l'encadré de droite est très serrée : « … d'ens.\
quasi-simplicial \ill{} \ill{} $\sigma$ tel que $D(\sigma) \neq
\emptyset$ » ; à gauche, reliée par une flèche : « (\emph{NB} $I_0 =
\delta(\Phi)$) »}

Le cas qui nous intéressera plus particulièrement est le cas
\[
  I = \mathbf{N}
\]
auquel cas $\delta$ prend le nom de fonction dimension, \struck{\ill{}} et
$\sup_{\overline{\mathbf{R}}} \delta(I)$ \struck{\ill{}} \add{s'appelle}
la dimension de la géom.\ de drapeaux $D$ (\uncertain{munie de sa} \struck{\ill{}}
fonction dimension) si elle est finie, mais il y a souvent \uncertain{intérêt} à remplacer
$\mathbf{N}$ par $[0,n]_{\mathbf{N}}$…
\note{« $\sup_{\overline{\mathbf{R}}}$ » : sous « sup », un
$\overline{\mathbf{R}}$ ; le mot initial de la ligne est biffé}

\note{un grand crochet vertical sépare ce qui précède de la suite}

Plus généralement, \add{soit une géom.\ de drapeaux $D$ donnée} supposons
donnée une loi qui associe à chaque \struck{partie} simplexe $\sigma
\subset \Phi$ (i.e.\ $\sigma \in \mathfrak{P}_f^*(\Phi)$ de la forme $I_d$)
une \emph{ordre total} sur $\sigma$, de façon que ces ordres totaux
s'induisent mutuellement [ce qui est le cas pour les ordres totaux induits
par un $\delta : \Phi \to I$ si $I$ lui-même est tot.\ ordonné, ou mieux
d'un \add{\uncertain{ouvert}} \struck{\ill{}} ordre \struck{\ill{}} sur
\uncertain{I} tel que les simplexes de $I$ soient tot.\ ord…) mais cet
\uncertain{ordre} : prendre sur $I$ \uncertain{l'ordre} \add{\ill{}}
\ill{} … ] Alors, posant pour $n \in \mathbf{N}$ ($n \geqslant 0$)
\begin{align*}
  D[n] &= \lbrace d \in D \mid \mathrm{card}\, I_d = n+1 \rbrace \\
  &= \coprod_{\sigma \in \mathfrak{P}_{n+1}(\Phi)} D(\sigma) \\
  & \Big( \simeq \text{aussi } \coprod_{\sigma \in
  \mathfrak{P}_{n+1}(I)} D'(\sigma) \Big)
\end{align*}
\note{« aussi » est une lecture \uncertain{incertaine}}

\page{18}
On trouve que $n \mapsto D[n]$ devient un ens.\ quasi-simplicial, i.e.\
un \add{contra}foncteur sur la catégorie $\Delta$ des ens.\ finis non vides
de la forme $\Delta_n = [0,n]_{\mathbf{N}}$, avec comme morphismes les
appl.\ \emph{injectives}. Cet ens.\ quasi-simplicial satisfait la
condition que
\[
  \forall\, n, m \in \mathbf{N} \text{ avec } n < m, \text{ et } d_m \in
  D[m], \text{ l'application}
\]
\[
  \mathrm{Hom}_\Delta(\Delta_n, \Delta_m) \to D_n, \qquad u \mapsto
  u^*(d_m)
\]
\marginal{\ill{}}
est injective, qui signifie aussi que la catégorie
\[
  \widehat{\Delta}/D[\,]
\]
est ordonnée (i.e.\ \uncertain{équivalente}, \uncertain{pardonnez}…).
L'ens.\ ordonné $D$ se reconstitue en fonction de $D[\,]$, comme l'ens.\
ordonné \uncertain{grandeur} par ladite catégorie ordonnée.
\struck{\ill{}} \add{\uncertain{On trouve}}
\struck{Le \ill{} de l'ens.}

\[
  \begin{array}{l}
    \text{géométrie des drapeaux } D \\
    \text{avec syst.\ d'ordres totaux sur} \\
    \text{les simplexes de } \Phi \text{ (ens.\ des} \\
    \text{él.\ minimaux de } D)
  \end{array}
  \Longleftrightarrow
  \begin{array}{l}
    \text{ens.\ quasi-} \\
    \text{simpliciaux} \\
    \text{satisfaisant la} \\
    \text{condition dess.}
  \end{array}
\]
\note{« totaux » est ajouté au-dessus de la ligne ; « avec » et « ens. »
sont des lectures \uncertain{incertaines}}

\page{19}
Soit $D$ une géométrie de drapeaux quasi-simpliciale
\struck{$D : \mathfrak{P}_f^*(\mathbf{N})^\circ \to (\mathrm{Ens})$}
\struck{(« géométrie de drapeaux avec fonction dimension »)}
\note{ces deux lignes sont barrées d'un zigzag}
[i.e.\ avec \add{$D$} un syst.\ d'ordres totaux compatibles sur les
simplexes de $\Phi = D_0$ (ens.\ des él.\ minimaux de $\Phi$)]
\note{il écrit bien « de $\Phi$ » à la fin, pour « de $D$ »}

d'où
\[
  X = |D|
\]
espace topologique, muni de parties fermées (compactes, des simplexes)
\[
  X(d) \qquad d \in D
\]
définissant une \emph{stratification} \uncertain{canonique} de $X$ par des simplexes
(les $X^\circ(d) = X(d) - \bigcup_{d' < d} X(d')$ forment une partition
de $X$ par des simplexes ouverts…)

Pour tout $F \in \Phi$, soit
\[
  \left\lbrace
  \begin{array}{l}
    |F| = \displaystyle\bigcup_{\substack{d \in D_* \\ \text{t.q.\ }
    \sup(I_d) = F}} X(d) \\[3ex]
    |F|^\circ = \displaystyle\bigcup_{\text{idem}} X(d)^\circ
  \end{array}
  \right.
\]
\note{à droite, relié par un long trait courbe : « $|F| - |F|^\circ
\overset{?}{=} \bigcup_{F' < F} |F'|$ », et au-dessous « $= \bigcup_{d \in
D_*,\ \sup I_d < F} X(d)$ », avec au-dessus « $\bigcup_{F' < F}
|F'|^\circ$ » et, en bas, « donc $(|F|)_{F \in \Phi}$ »}

\marginal{?}
Quand peut-on dire que la famille des \add{et celle des cellules fermées}
\struck{\ill{}} forment celle des cellules \emph{fermées} d'une
\emph{stratification} de $X$, dont les cellules ouvertes sont les
$|F|^\circ$ ?

\add{En fait :}
\struck{a) $(|F|)_{F \in \Phi}$ \ill{} $|F| = \bigcup X(d) = X$}
\add{loc.\ fermés ??}
\begin{itemize}
  \item[a)] Les $|F|^\circ$ forment \emph{partition} de $X$ [\emph{NB}
  $|F|^\circ \supseteq$ \struck{\ill{}} \add{les \uncertain{sommets}
  \uncertain{maximaux}} donc $|F|^\circ \neq \emptyset$ ]
  \item[b)] $|F| = \overline{|F|^\circ}$ car $X(d) = \overline{X(d)^\circ}$
  \quad ) \uncertain{\ill{} de « cellules » ouvertes}
  \item[c)] \struck{$|F|^\circ = |F| - \bigcup (X(d))$, $d \in D_*$,
  $\sup(I_d) = F$}
\end{itemize}

\page{20}
Soit $\Delta$ la catégorie des ens.\ finis tot\add{alement} ordonnés
\add{non vides}, avec comme morphismes les applications
\emph{strictement} croissantes. Soit
\[
  D_* \in \mathrm{Ob}\, \widehat{\Delta} \qquad D_* : \Delta^\circ
  \longrightarrow (\mathrm{Ens})
\]
donc $D_*$ définit une catégorie fibrée : fibres discrètes sur $\Delta$,
à quelle condition celle-ci est-elle ordonnée \add{(ou préordonnée,
\uncertain{cela revient au même})} ? Posons $\Delta_n = [0,n] \in
\mathrm{Ob}\,\Delta$, $D_n = D_*(\Delta_n)$, la condition est celle-ci

(1) $\forall\, n, m \in \mathbf{N}$, avec $n < m$, et $d_m \in D_m$,
l'application
\[
  \mathrm{Hom}_\Delta(\Delta_n, \Delta_m) \longrightarrow D_n, \qquad u
  \longmapsto u^*(d_m)
\]
est injective.
\note{le « 1 » est cerclé}

Ceci posé, on a donc un ens.\ ordonné $D_* = \coprod_{m \in \mathbf{N}}
D_m$ et un foncteur fibrant
\[
  \underline{D}_* \longrightarrow \Delta
\]
Considérons la réalisation topologique $|D_*|$ de $D_*$, donc $\forall\, d
\in D_* = \mathrm{Ob}\,\underline{D}$, on a une partie \add{fermée} $|d|
\subset |D_*|$, l'application $d \mapsto |d|$ de
\[
  D_* \to \mathfrak{P}(|D_*|)
\]
\add{\uncertain{est}} \uncertain{injective} \add{\uncertain{et}}
\uncertain{respecte} la relation d'ordre. Considérons \add{pour $d \in
D_m$} \struck{\ill{}} l'ens.\ $D_d = \lbrace \delta \in D \mid \delta
\leqslant d \rbrace$, on trouve qu'il est canoniquement isomorphe (comme
ens.\ ordonné) à $\mathfrak{P}^*(\Delta_m)$ ($\mathfrak{P}^*$ désigne
l'ens.\ des parties $\neq \emptyset$). Donc $|d|$ est canoniquement
isomorphe au simplexe type $|\Delta_m|$. Ainsi on trouve
\[
  |D_*| = \varinjlim_{d \in D_*} |d| = \varinjlim_{d \in D_*} |\Delta_m| =
  \varinjlim_{\Delta \in \Delta/D} |\Delta|
\]

\end{document}
